% A.78 reserved by the root.
\subsection{Exact ordinary bad-place values and the remaining cusp exponents}

\begin{statement}{A.78}{Proved}
{The unknown exponent vector is supported only at actual cuspidal places}
Let $s$ be a noncuspidal Selmer section in A.14 with the
labeled integral realization $E_{\lambda_0}$ of A.16.
At every potentially multiplicative prime $p$ where $s_p$
is the section of an ordinary point $x_p$, one has the exact
equality
\[
                              x_p=\lambda_0.
\tag{A.78.1}
\]
This assertion requires the given labeled integral comparison,
with no condition on its Weil multiplier.

Now suppose that the full Legendre congruence images of $s$
and $s_{\lambda_0}$ are conjugate by $H$.
Write
$\lambda_0=a/c$, $1-\lambda_0=b/c$, with
$a+b=c$, $c>0$, and $a,b,c$ pairwise coprime nonzero integers.
Let $\mathcal C$ be the finite set of primes at which $s_p$
is cuspidal.  At each $p\in\mathcal C$, its label is the
contact label of $\lambda_0$:
\[
 c(p)=
 \begin{cases}
 0,&p\mid a,\\
 1,&p\mid b,\\
 \infty,&p\mid c.
 \end{cases}
 \qquad
 e_p=v_p(abc),\quad
 m_p=e_p-4\,\mathbf1_{p=2}>0.
\tag{A.78.2}
\]
Put $\mathcal C_i=\{p\in\mathcal C:c(p)=i\}$.
There are exponents $\delta_p\in24\widehat\Z$ with
\[
 \begin{split}
 A&=\lambda_0
          \prod_{p\in\mathcal C_0\cup\mathcal C_\infty}
                                                        p^{\delta_p},\\
 B&=(1-\lambda_0)
          \prod_{p\in\mathcal C_1\cup\mathcal C_\infty}
                                                        p^{\delta_p}.
 \end{split}
\tag{A.78.3}
\]
The signs in these expressions are the signs of the displayed
rational factors.  More precisely, for each $p\in\mathcal C$
there is a unit $u_p\in\widehat\Z^\times$ such that
\[
 \delta_p=\sigma_p m_p(u_p^2-1),
 \qquad
 \sigma_p=
 \begin{cases}1,&c(p)=0\text{ or }1,\\-1,&c(p)=\infty.
 \end{cases}
\tag{A.78.4}
\]
The units $u_p$ are allowed to depend on the local prime $p$.
The assertion does not identify these units with one common
unit or force the exponents to vanish.
\end{statement}
\begin{proof}
We first prove (A.78.1), retaining the full integral extensions.
Choose a finite extension $F/\Q_p$ over which $E_{x_p}$ and
$E_{\lambda_0}$ are split Tate curves, with parameters
$q_x,q_0\in F^\times$ and normalized positive valuations
$N_x,N_0$.  The first curve is potentially multiplicative
because its integral system is that of the second; alternatively
its nonzero infinite Tate inertia is carried across the given
isomorphism, as in A.34.

For every coefficient prime $\ell$, including $p$, the
cyclotomic line in the rational Tate representation is unique
and its integral intersection is saturated.  The comparison
therefore has, in integral Tate bases, triangular matrix with
unit diagonal entries $a_\ell,d_\ell$ and satisfies
\[
 \kappa_\ell(q_x)=(a_\ell/d_\ell)\kappa_\ell(q_0),
 \qquad
                 N_x=(a_\ell/d_\ell)N_0.
\tag{A.78.5}
\]
These are the actual extension and valuation maps of A.34.
In particular the positive rational number $N_x/N_0$ is a
unit at every rational prime.  It equals one, so
$N_x=N_0$ and $a_\ell=d_\ell$ for every $\ell$.
Consequently $\kappa_\ell(q_x/q_0)=0$ for every coefficient.
The map $F^\times\to\widehat{F^\times}$ is injective:
its decomposition into a uniformizer factor, finite roots of
unity and a principal-unit group reduces this assertion to
$\Z\hookrightarrow\widehat\Z$, to a finite group, and to a
finite-rank $\Z_p$-module.  Thus
\[
                                  q_x=q_0.
\tag{A.78.6}
\]
The actual Tate uniformizations now give an elliptic-curve
isomorphism $\phi:E_{\lambda_0,F}\to E_{x_p,F}$.

To recover the labels, compare $T_2(\phi)$ with the given
integral $2$-adic intertwiner.  Their quotient commutes with
the Tate representation.  Its rational matrix is triangular,
say $\left(\begin{smallmatrix}a&b\\0&d\end{smallmatrix}\right)$.
For a cocycle $c$ representing $\kappa_2(q_0)$, commutation
implies
\[
                    (a-d)c(g)=(\chi_2(g)-1)b.
\tag{A.78.7}
\]
The extension class has nonzero valuation $N_0$, hence is
nonzero over $\Q_2$.  Taking cohomology in (A.78.7) gives
$a=d$; since $\chi_2$ is nontrivial on $G_F$, the same
equation then gives $b=0$.
The integral quotient is therefore a scalar in $\Z_2^\times$,
whose reduction modulo two is one.
The original comparison preserves the ordered nonzero
$2$-torsion points, so $\phi$ preserves their labels as well.
An origin-preserving isomorphism between these characteristic-zero
Weierstrass equations has coordinate map
$X\mapsto u^2X+r$, $Y\mapsto u^3Y$.
Preserving $(0,0)$ and $(1,0)$ forces $r=0$ and $u^2=1$;
preserving the third point then gives $x_p=\lambda_0$.
This proves (A.78.1) without a determinant restriction.

We next fix the degeneration label, including at two.
For odd $p$ the saturated cyclotomic $T_2$ line is the line
of the noncolliding labeled $2$-torsion point, by the exact
calculation in A.41.  Its values for contact labels
$0,1,\infty$ are respectively
$\langle(1,0)\rangle$, $\langle(0,0)\rangle$ and
$\langle(\lambda_0,0)\rangle$.

Here is an explicit valuation calculation proving the same
identification in residue characteristic two.  Over a further
finite extension write $q=r^2$, with $v(r)>0$, and put
$f(z)=z/(1-z)^2$.  The Tate coordinate is
\[
 X(z)=\sum_{n\in\Z}f(q^nz)
                      -2\sum_{n\ge1}\frac{nq^n}{1-q^n},
\tag{A.78.8}
\]
the convergent uniformization formula of
\cite[equations (1)--(4) and Theorem 1]{LegTate1995}.
The constant term in Tate's equation (1) has this form
because the coefficient of $q^k$ in each of
$\sum_{n\ge1}q^n/(1-q^n)^2$ and
$\sum_{n\ge1}nq^n/(1-q^n)$ is $\sum_{d\mid k}d$.
The identity $f(z)=f(z^{-1})$ pairs the terms for $z=\pm r$.
Writing $d=v(2)$ and $h=v(r)>0$ gives
\[
 \begin{split}
 X(-1)&=-\tfrac14+\text{terms of positive valuation},\\
 X(r)-X(-r)
   &=4\sum_{j\ge0}
       \frac{r^{2j+1}(1+r^{4j+2})}{(1-r^{4j+2})^2},\\
 v(X(-1))&=-2d,\qquad
 v(X(r)-X(-r))=2d+h,\\
 v(X(\pm r)-X(-1))&=-2d.
 \end{split}
\tag{A.78.9}
\]
For the second valuation the first summand has strictly smaller
valuation than every later summand; each $X(\pm r)$ itself
has positive valuation.  Thus the uniquely closest pair of
Tate $2$-torsion coordinates is the pair represented by
$\pm r$, while $-1$ represents the cyclotomic point.
An elliptic-curve isomorphism rescales all differences of
$X$ coordinates by the same nonzero square.  It therefore
preserves this uniquely closest pair.
In the Legendre equation, contact at zero makes
$0,\lambda_0$ the closest pair; contact at one makes
$1,\lambda_0$ the closest pair; and contact at infinity
makes $0,1$ the closest pair after a common scaling.
The last identification also follows from the actual map
$X=tX'$, $Y=t\sqrt t\,Y'$ to parameter $1/t$.
This proves precisely the three labeled lines just stated
also at two.

On a cusp packet the same lines are computed on the formal
Tate curve over the characteristic-zero field $\Q_p((w))$,
with $w=t,1-t,1/t$ respectively.  Pullback preserves its
cyclotomic submodule.  Its extension cannot split under the
comparison with the realized curve, because the latter has
nonzero rational Tate class.  Thus its saturated line is
carried to that of the realized curve.  The comparison in
A.16 is the identity on the ordered $2$-torsion, so the three
distinct lines prove (A.78.2)'s label assertion at every prime.
The exact formula
$v_p(j(E_{\lambda_0}))=8v_p(2)-2e_p$
at a contact prime shows that its Tate parameter has valuation
$2m_p>0$.  A.37 already confines every local cusp to such a
potentially multiplicative prime, so $\mathcal C$ is finite.

We now calculate the complete global support.
At a potentially multiplicative prime outside $\mathcal C$,
(A.78.1) fixes both ordinary local valuations exactly.
At every potentially good odd prime, the Legendre $j$ formula
forces an ordinary point to have both coordinates units;
otherwise its $j$ valuation would be negative.
At two, a potentially good ordinary point has exactly one
of the twelve valuation pairs
\[
       (k,0),\quad(0,k),\quad(-k,-k),\qquad1\le k\le4.
\tag{A.78.10}
\]
Indeed, two odd units cannot sum to one; the ultrametric
inequality gives these three forms, and the $j$ valuation
$8-2k\ge0$ gives $k\le4$.
The twelve pairs have distinct reductions modulo eight.
The full-$H$ comparison, under the exact character map of A.56,
gives
\[
 A/\lambda_0=1,\quad B/(1-\lambda_0)=1
               \quad\text{modulo eighth powers},
 \qquad
 \frac{AB}{\lambda_0(1-\lambda_0)}=1
               \quad\text{modulo cubes}.
\tag{A.78.11}
\]
Applying the ordinary valuation map at two therefore fixes
the pair in (A.78.10) to be that of $\lambda_0$.
This treats every prime outside $\mathcal C$.

The sign and all prime valuations inject into
$\widehat{\Q^\times}$, as in A.39.
The first two classes in (A.78.11) have positive sign, since
the sign is detected modulo eighth powers.
They consequently have product expansions supported only
on $\mathcal C$.  At a cusp of label zero the second
coordinate equals one locally, so its valuation is zero;
at a cusp of label one the first valuation is zero; and
at infinity the two valuations are equal because $B/A=-1$.
Since the realized rational point has the same contact label,
the difference of valuation pairs is respectively
$(\delta_p,0)$, $(0,\delta_p)$ or $(\delta_p,\delta_p)$.
This proves the exact product expansions (A.78.3).

For completeness, the map (A.78.11) sends a relative
valuation pair $(\alpha,\beta)$ at any prime to
\[
              (\alpha\bmod8,\beta\bmod8,
                                      \alpha+\beta\bmod3).
\tag{A.78.12}
\]
Its kernel is the closed subgroup generated by
$(24,0),(0,24),(8,-8)$ in A.56.
On any of the three just identified label vectors, lying
in this kernel is exactly $\delta_p\in24\widehat\Z$.
For the infinity vector one uses that multiplication by
two is invertible modulo three.

Finally retain the determinant in the local comparison.
After a common finite splitting extension, the global
$H$ conjugator and paired Tate bases give a triangular
matrix with diagonal entries $a_{p,\ell},d_{p,\ell}$
satisfying $a_{p,\ell}d_{p,\ell}=1$.
The exact cusp Tate classes are
$A^2/256$, $B^2/256$ and $A^{-2}/256$, with the
quadratic twists and formal-unit specialization of A.63.
Their valuation comparison with the realized Tate class,
after canceling the common ramification index, is
\[
          2m_p+2\sigma_p\delta_{p,\ell}
                           =2m_p a_{p,\ell}^{2}.
\tag{A.78.13}
\]
All diagonal entries are integral units because the
cyclotomic lines are saturated.  Equation (A.78.13)
in every torsion-free $\Z_\ell$ gives (A.78.4) with
$u_p=(a_{p,\ell})_\ell$.
This uses the full extensions and their Weil pairings;
it does not replace the determinant-one condition by
an unspecified linear intertwiner.
\end{proof}

\paragraph{Exact simultaneous equations on the remaining cusp set.}
The regular coordinate at each local cusp yields the following
equalities in $\widehat{\Q_p^\times}$, with the products
localized through the ordinary rational prime inclusions:
\[
\begin{array}{c|c}
 p\in\mathcal C_0&
 (1-\lambda_0)\displaystyle\prod_{r\in\mathcal C_1\cup\mathcal C_\infty}
                                      r^{\delta_r}=1\\[2mm]
 p\in\mathcal C_1&
 \lambda_0\displaystyle\prod_{r\in\mathcal C_0\cup\mathcal C_\infty}
                                      r^{\delta_r}=1\\[2mm]
 p\in\mathcal C_\infty&
 \dfrac{\lambda_0-1}{\lambda_0}
       \displaystyle\prod_{r\in\mathcal C_1}r^{\delta_r}
       \prod_{r\in\mathcal C_0}r^{-\delta_r}=1.
\end{array}
\tag{A.78.14}
\]
These equalities retain the torsion components as well as
the coefficient-prime logarithms.
For example, for $p\in\mathcal C_0$ their logarithmic
consequence, after substituting (A.78.4), is the exact
unit-variable quadratic equation
\[
 \begin{aligned}
 &\sum_{r\in\mathcal C_1}m_r\log_p(r)u_{r,p}^2
 -\sum_{r\in\mathcal C_\infty}m_r\log_p(r)u_{r,p}^2\\
 &\hspace{8mm}=-\log_p(1-\lambda_0)
   +\sum_{r\in\mathcal C_1}m_r\log_p(r)
   -\sum_{r\in\mathcal C_\infty}m_r\log_p(r),
 \qquad u_{r,p}\in\Z_p^\times.
 \end{aligned}
\tag{A.78.15}
\]
Here every $r$ in a sum differs from $p$, so its logarithm
is the usual logarithm of a local unit, killing its finite
root-of-unity factor.  The other two rows give corresponding
equations by the same substitution into (A.78.14).
Thus each proposed cusp set has an explicit finite system
of integral unit-square conditions.  A contradiction modulo
a finite power of a local prime excludes that set, since it
is the image of these actual equalities.  The full Tate-class
equations of A.63 remain in addition to these regular-coordinate
equations.  No assertion that all such systems are inconsistent
is made here.

\noindent The finite certificate
\path{code/legendre_ordinary_support_certificate.py}
and its retained output check the paired dyadic Tate series,
their closest-pair valuations and resulting Legendre $j$
identity, as well as the finite sign/valuation character maps.
The source \cite{LegTate1995} was opened directly for its
explicit formulas and uniformization theorem; the infinite
valuation and Galois assertions are proved above.
